\documentclass[11pt,a4paper]{article}
\usepackage[utf8]{inputenc}
\usepackage[T1]{fontenc}
\usepackage[spanish,es-minimal,es-noshorthands]{babel}
\usepackage{lmodern}
\usepackage[sfdefault,lining]{FiraSans}
\usepackage[a4paper,top=2.3cm,bottom=2.2cm,left=1.8cm,right=1.8cm,headheight=1.3cm,headsep=0.4cm,footskip=1cm]{geometry}
\usepackage{amsmath,amssymb}
\allowdisplaybreaks
\usepackage{icomma}
\usepackage{graphicx}
\usepackage[table]{xcolor}
\usepackage{tikz}
\usepackage{fontawesome5}
\usepackage{ccicons}
\usepackage{fancyhdr}
\usepackage{lastpage}
\usepackage{tasks}
\usepackage{needspace}
\usepackage[most]{tcolorbox}
\usepackage{array}
\usepackage[hidelinks]{hyperref}
\AtBeginDocument{\hypersetup{pdftitle={Ficha de ejercicios de la unidad 8: Representación gráfica de funciones. 2º Bach.},pdfsubject={Matemáticas II 2º Bach.},pdfkeywords={representación de funciones, dominio, simetrías, periodicidad, asíntotas, ramas parabólicas, monotonía, curvatura, funciones radicales, trigonométricas, a trozos, transformaciones, PAU}}}

\definecolor{navy}{HTML}{1B2748}
\definecolor{accent}{HTML}{E8603A}
\definecolor{teal}{HTML}{1F9E90}
\definecolor{soft}{HTML}{EEF1F8}
\definecolor{muted}{HTML}{5D6685}

\setlength{\parindent}{0pt}
\setlength{\parskip}{3pt}
\settasks{label=\alph*), label-format=\bfseries\color{navy}, label-width=1.3em, item-indent=1.8em, column-sep=0.6em, after-item-skip=5pt, before-skip=3pt, after-skip=4pt}

% ---------- capçalera i peu ----------
\newcommand{\llicencia}{\href{https://creativecommons.org/licenses/by-nc-sa/4.0/deed.es}{\ccbyncsa}\ \ Licencia Creative Commons Atribución-NoComercial-CompartirIgual 4.0 Internacional (CC BY-NC-SA 4.0)}
\pagestyle{fancy}
\fancyhf{}
\fancyhead[L]{\raisebox{-0.35\height}{\includegraphics[height=0.95cm]{logo-gva.png}}}
\fancyhead[C]{\small\color{muted}Ficha de ejercicios · Unidad 8: Representación gráfica de funciones · 2\textsuperscript{o} Bach.}
\fancyhead[R]{\raisebox{-0.35\height}{\includegraphics[height=0.95cm]{logo-vermellar.png}}}
\fancyfoot[L]{\scriptsize\color{muted}\llicencia}
\fancyfoot[R]{\small\color{muted}\thepage}
\renewcommand{\headrulewidth}{0.4pt}
\renewcommand{\headrule}{\color{muted!40}\hrule height \headrulewidth}
\fancypagestyle{primera}{\fancyhf{}\renewcommand{\headrulewidth}{0pt}
  \fancyfoot[L]{\scriptsize\color{muted}\llicencia}
  \fancyfoot[R]{\small\color{muted}\thepage}}

% ---------- nivells i exercicis ----------
\newcommand{\estrelles}[1]{{\color{accent}\ifcase#1\or\faStar\faStar[regular]\faStar[regular]\or\faStar\faStar\faStar[regular]\or\faStar\faStar\faStar\fi}}
\newcounter{exer}
\newcommand{\ex}[1]{\par\needspace{4\baselineskip}\medskip\refstepcounter{exer}\noindent{\color{navy}\bfseries\theexer.}\ \ }

\newcommand{\seccio}[2]{\par\needspace{10\baselineskip}\bigskip
  \begin{tcolorbox}[colback=navy,colframe=navy,coltext=white,arc=2mm,boxrule=0pt,left=3mm,right=3mm,top=1.2mm,bottom=1.2mm]
  {\large\bfseries #1\quad #2}
  \end{tcolorbox}\vspace{-2pt}}


\setlength{\emergencystretch}{3em}

% quadre de recordatori (només a la fitxa d'exercicis)
\newcommand{\recorda}[1]{\par\smallskip\begin{tcolorbox}[colback=accent!7,colframe=accent!55,arc=2mm,boxrule=0.5pt,left=2.5mm,right=2.5mm,top=1mm,bottom=1mm]\small{\color{accent}\faLightbulb[regular]}\ \textbf{Recuerda:}\ #1\end{tcolorbox}}
% etiqueta dels exercicis de la PAU
\newcommand{\pau}[1]{{\color{teal}\small\bfseries\faUniversity\ PAU · #1.}\ }
\newcommand{\res}[1]{#1}\newcommand{\rest}[1]{#1}

% =====================================================================
% Capçalera FentMates de les fitxes: un color per a cada curs.
% Els colors són els de la pàgina inicial (LEVELS a assets/js/materials.js);
% el logo ve de _eines/logo/pdf/ (python3 _eines/logo/genera.py).
%
% Ús, al final del preàmbul de la unitat:
%   % =====================================================================
% Capçalera FentMates de les fitxes: un color per a cada curs.
% Els colors són els de la pàgina inicial (LEVELS a assets/js/materials.js);
% el logo ve de _eines/logo/pdf/ (python3 _eines/logo/genera.py).
%
% Ús, al final del preàmbul de la unitat:
%   % =====================================================================
% Capçalera FentMates de les fitxes: un color per a cada curs.
% Els colors són els de la pàgina inicial (LEVELS a assets/js/materials.js);
% el logo ve de _eines/logo/pdf/ (python3 _eines/logo/genera.py).
%
% Ús, al final del preàmbul de la unitat:
%   \input{../comu/fentmates.tex}
%   \fmcurs{1eso}
% i, al començament del document:
%   \capcalera{Matemàtiques}{1r ESO}{Fitxa d'exercicis · Unitat 2}{Divisibilitat}
% Opcional abans de \input: \def\fmneed{14} (espai mínim abans de cada secció).
% =====================================================================
\makeatletter
\def\fm@logodir{../../logo/pdf}
\providecommand{\fmneed}{10}
\def\fm@cursnom{}

% clau del curs, color i color del text damunt del color
\newcommand{\fm@curs}[3]{\@namedef{fm@color@#1}{#2}\@namedef{fm@text@#1}{#3}}
\fm@curs{1eso}{E6B84A}{navy}
\fm@curs{2eso}{3B82F6}{white}
\fm@curs{3eso}{FF7A50}{white}
\fm@curs{4eso-a}{A78BFA}{white}
\fm@curs{4eso-b}{6D28D9}{white}
\fm@curs{1batx-mi}{9B1C1C}{white}
\fm@curs{1batx-ccss}{EF6B6B}{white}
\fm@curs{2batx-mii}{1E6B47}{white}
\fm@curs{2batx-ccss}{2BB3A3}{white}

\newcommand{\fmcurs}[1]{%
  \@ifundefined{fm@color@#1}{\PackageError{fentmates}{Curs desconegut: #1}{}}{}%
  \edef\fm@tmp{\noexpand\definecolor{curs}{HTML}{\@nameuse{fm@color@#1}}%
               \noexpand\colorlet{curstext}{\@nameuse{fm@text@#1}}}\fm@tmp
  \colorlet{cursfosc}{curs!72!black}%
  \def\fm@logo{\fm@logodir/vertical-#1.pdf}%
  \def\fm@marca{\fm@logodir/marca-#1.pdf}}

% ---------- capçalera de la primera pàgina ----------
% #1 matèria (petit), #2 curs, #3 tipus de document i unitat, #4 títol
\newcommand{\capcalera}[4]{%
  \gdef\fm@cursnom{#2}%
  \par\noindent
  \begin{tcolorbox}[colback=curs,colframe=curs,boxrule=0pt,arc=2.2mm,width=3.5cm,height=2.15cm,
      valign=center,left=3mm,right=2mm,top=1mm,bottom=1mm,nobeforeafter,box align=center]
    \color{curstext}{\footnotesize #1}\par\vspace{1pt}{\fontsize{19}{21}\selectfont\bfseries #2\par}%
  \end{tcolorbox}\hspace{5mm}%
  \begin{minipage}[c]{\dimexpr\linewidth-3.5cm-5mm-2.9cm\relax}
    {\color{cursfosc}\small\bfseries\MakeUppercase{#3}}\par\vspace{2pt}
    {\color{navy}\fontsize{20}{23}\selectfont\bfseries #4\par}
  \end{minipage}\hfill
  \begin{minipage}[c]{2.6cm}\raggedleft\includegraphics[width=2.6cm]{\fm@logo}\end{minipage}\par
  \vspace{2pt}}

% ---------- barres de secció del color del curs ----------
\renewcommand{\seccio}[2]{\par\needspace{\fmneed\baselineskip}\bigskip
  \begin{tcolorbox}[colback=curs,colframe=curs,coltext=curstext,arc=2mm,boxrule=0pt,left=3mm,right=3mm,top=1.2mm,bottom=1.2mm]
  {\large\bfseries #1\quad #2}
  \end{tcolorbox}\vspace{-2pt}}

% ---------- estrelles de dificultat del color del curs (a Batxillerat no se'n posen) ----------
\@ifundefined{estrelles}{}{%
  \renewcommand{\estrelles}[1]{{\color{curs}\ifcase#1\or\faStar\faStar[regular]\faStar[regular]\or\faStar\faStar\faStar[regular]\or\faStar\faStar\faStar\fi}}}

% ---------- capçalera i peu de les pàgines ----------
% primera pàgina: logos institucionals al peu
\fancypagestyle{primera}{\fancyhf{}\renewcommand{\headrulewidth}{0pt}
  \fancyfoot[L]{\raisebox{-0.32\height}{\includegraphics[height=0.78cm]{logo-gva.png}}\hspace{5mm}%
    \raisebox{-0.32\height}{\includegraphics[height=0.78cm]{logo-vermellar.png}}}
  \fancyfoot[R]{\scriptsize\color{muted}\href{https://creativecommons.org/licenses/by-nc-sa/4.0/deed.ca}{\ccbyncsa}\ \ CC BY-NC-SA 4.0 · fentmates.org\qquad\small\thepage}}
% la resta: el símbol a l'esquerra i el curs a la dreta, amb el títol al centre com abans
\fancyhead[L]{\raisebox{-0.3\height}{\includegraphics[height=0.62cm]{\fm@marca}}}
\fancyhead[R]{\small\bfseries\color{cursfosc}\fm@cursnom}
\renewcommand{\headrule}{\color{curs!45}\hrule height \headrulewidth}
\makeatother

%   \fmcurs{1eso}
% i, al començament del document:
%   \capcalera{Matemàtiques}{1r ESO}{Fitxa d'exercicis · Unitat 2}{Divisibilitat}
% Opcional abans de \input: \def\fmneed{14} (espai mínim abans de cada secció).
% =====================================================================
\makeatletter
\def\fm@logodir{../../logo/pdf}
\providecommand{\fmneed}{10}
\def\fm@cursnom{}

% clau del curs, color i color del text damunt del color
\newcommand{\fm@curs}[3]{\@namedef{fm@color@#1}{#2}\@namedef{fm@text@#1}{#3}}
\fm@curs{1eso}{E6B84A}{navy}
\fm@curs{2eso}{3B82F6}{white}
\fm@curs{3eso}{FF7A50}{white}
\fm@curs{4eso-a}{A78BFA}{white}
\fm@curs{4eso-b}{6D28D9}{white}
\fm@curs{1batx-mi}{9B1C1C}{white}
\fm@curs{1batx-ccss}{EF6B6B}{white}
\fm@curs{2batx-mii}{1E6B47}{white}
\fm@curs{2batx-ccss}{2BB3A3}{white}

\newcommand{\fmcurs}[1]{%
  \@ifundefined{fm@color@#1}{\PackageError{fentmates}{Curs desconegut: #1}{}}{}%
  \edef\fm@tmp{\noexpand\definecolor{curs}{HTML}{\@nameuse{fm@color@#1}}%
               \noexpand\colorlet{curstext}{\@nameuse{fm@text@#1}}}\fm@tmp
  \colorlet{cursfosc}{curs!72!black}%
  \def\fm@logo{\fm@logodir/vertical-#1.pdf}%
  \def\fm@marca{\fm@logodir/marca-#1.pdf}}

% ---------- capçalera de la primera pàgina ----------
% #1 matèria (petit), #2 curs, #3 tipus de document i unitat, #4 títol
\newcommand{\capcalera}[4]{%
  \gdef\fm@cursnom{#2}%
  \par\noindent
  \begin{tcolorbox}[colback=curs,colframe=curs,boxrule=0pt,arc=2.2mm,width=3.5cm,height=2.15cm,
      valign=center,left=3mm,right=2mm,top=1mm,bottom=1mm,nobeforeafter,box align=center]
    \color{curstext}{\footnotesize #1}\par\vspace{1pt}{\fontsize{19}{21}\selectfont\bfseries #2\par}%
  \end{tcolorbox}\hspace{5mm}%
  \begin{minipage}[c]{\dimexpr\linewidth-3.5cm-5mm-2.9cm\relax}
    {\color{cursfosc}\small\bfseries\MakeUppercase{#3}}\par\vspace{2pt}
    {\color{navy}\fontsize{20}{23}\selectfont\bfseries #4\par}
  \end{minipage}\hfill
  \begin{minipage}[c]{2.6cm}\raggedleft\includegraphics[width=2.6cm]{\fm@logo}\end{minipage}\par
  \vspace{2pt}}

% ---------- barres de secció del color del curs ----------
\renewcommand{\seccio}[2]{\par\needspace{\fmneed\baselineskip}\bigskip
  \begin{tcolorbox}[colback=curs,colframe=curs,coltext=curstext,arc=2mm,boxrule=0pt,left=3mm,right=3mm,top=1.2mm,bottom=1.2mm]
  {\large\bfseries #1\quad #2}
  \end{tcolorbox}\vspace{-2pt}}

% ---------- estrelles de dificultat del color del curs (a Batxillerat no se'n posen) ----------
\@ifundefined{estrelles}{}{%
  \renewcommand{\estrelles}[1]{{\color{curs}\ifcase#1\or\faStar\faStar[regular]\faStar[regular]\or\faStar\faStar\faStar[regular]\or\faStar\faStar\faStar\fi}}}

% ---------- capçalera i peu de les pàgines ----------
% primera pàgina: logos institucionals al peu
\fancypagestyle{primera}{\fancyhf{}\renewcommand{\headrulewidth}{0pt}
  \fancyfoot[L]{\raisebox{-0.32\height}{\includegraphics[height=0.78cm]{logo-gva.png}}\hspace{5mm}%
    \raisebox{-0.32\height}{\includegraphics[height=0.78cm]{logo-vermellar.png}}}
  \fancyfoot[R]{\scriptsize\color{muted}\href{https://creativecommons.org/licenses/by-nc-sa/4.0/deed.ca}{\ccbyncsa}\ \ CC BY-NC-SA 4.0 · fentmates.org\qquad\small\thepage}}
% la resta: el símbol a l'esquerra i el curs a la dreta, amb el títol al centre com abans
\fancyhead[L]{\raisebox{-0.3\height}{\includegraphics[height=0.62cm]{\fm@marca}}}
\fancyhead[R]{\small\bfseries\color{cursfosc}\fm@cursnom}
\renewcommand{\headrule}{\color{curs!45}\hrule height \headrulewidth}
\makeatother

%   \fmcurs{1eso}
% i, al començament del document:
%   \capcalera{Matemàtiques}{1r ESO}{Fitxa d'exercicis · Unitat 2}{Divisibilitat}
% Opcional abans de \input: \def\fmneed{14} (espai mínim abans de cada secció).
% =====================================================================
\makeatletter
\def\fm@logodir{../../logo/pdf}
\providecommand{\fmneed}{10}
\def\fm@cursnom{}

% clau del curs, color i color del text damunt del color
\newcommand{\fm@curs}[3]{\@namedef{fm@color@#1}{#2}\@namedef{fm@text@#1}{#3}}
\fm@curs{1eso}{E6B84A}{navy}
\fm@curs{2eso}{3B82F6}{white}
\fm@curs{3eso}{FF7A50}{white}
\fm@curs{4eso-a}{A78BFA}{white}
\fm@curs{4eso-b}{6D28D9}{white}
\fm@curs{1batx-mi}{9B1C1C}{white}
\fm@curs{1batx-ccss}{EF6B6B}{white}
\fm@curs{2batx-mii}{1E6B47}{white}
\fm@curs{2batx-ccss}{2BB3A3}{white}

\newcommand{\fmcurs}[1]{%
  \@ifundefined{fm@color@#1}{\PackageError{fentmates}{Curs desconegut: #1}{}}{}%
  \edef\fm@tmp{\noexpand\definecolor{curs}{HTML}{\@nameuse{fm@color@#1}}%
               \noexpand\colorlet{curstext}{\@nameuse{fm@text@#1}}}\fm@tmp
  \colorlet{cursfosc}{curs!72!black}%
  \def\fm@logo{\fm@logodir/vertical-#1.pdf}%
  \def\fm@marca{\fm@logodir/marca-#1.pdf}}

% ---------- capçalera de la primera pàgina ----------
% #1 matèria (petit), #2 curs, #3 tipus de document i unitat, #4 títol
\newcommand{\capcalera}[4]{%
  \gdef\fm@cursnom{#2}%
  \par\noindent
  \begin{tcolorbox}[colback=curs,colframe=curs,boxrule=0pt,arc=2.2mm,width=3.5cm,height=2.15cm,
      valign=center,left=3mm,right=2mm,top=1mm,bottom=1mm,nobeforeafter,box align=center]
    \color{curstext}{\footnotesize #1}\par\vspace{1pt}{\fontsize{19}{21}\selectfont\bfseries #2\par}%
  \end{tcolorbox}\hspace{5mm}%
  \begin{minipage}[c]{\dimexpr\linewidth-3.5cm-5mm-2.9cm\relax}
    {\color{cursfosc}\small\bfseries\MakeUppercase{#3}}\par\vspace{2pt}
    {\color{navy}\fontsize{20}{23}\selectfont\bfseries #4\par}
  \end{minipage}\hfill
  \begin{minipage}[c]{2.6cm}\raggedleft\includegraphics[width=2.6cm]{\fm@logo}\end{minipage}\par
  \vspace{2pt}}

% ---------- barres de secció del color del curs ----------
\renewcommand{\seccio}[2]{\par\needspace{\fmneed\baselineskip}\bigskip
  \begin{tcolorbox}[colback=curs,colframe=curs,coltext=curstext,arc=2mm,boxrule=0pt,left=3mm,right=3mm,top=1.2mm,bottom=1.2mm]
  {\large\bfseries #1\quad #2}
  \end{tcolorbox}\vspace{-2pt}}

% ---------- estrelles de dificultat del color del curs (a Batxillerat no se'n posen) ----------
\@ifundefined{estrelles}{}{%
  \renewcommand{\estrelles}[1]{{\color{curs}\ifcase#1\or\faStar\faStar[regular]\faStar[regular]\or\faStar\faStar\faStar[regular]\or\faStar\faStar\faStar\fi}}}

% ---------- capçalera i peu de les pàgines ----------
% primera pàgina: logos institucionals al peu
\fancypagestyle{primera}{\fancyhf{}\renewcommand{\headrulewidth}{0pt}
  \fancyfoot[L]{\raisebox{-0.32\height}{\includegraphics[height=0.78cm]{logo-gva.png}}\hspace{5mm}%
    \raisebox{-0.32\height}{\includegraphics[height=0.78cm]{logo-vermellar.png}}}
  \fancyfoot[R]{\scriptsize\color{muted}\href{https://creativecommons.org/licenses/by-nc-sa/4.0/deed.ca}{\ccbyncsa}\ \ CC BY-NC-SA 4.0 · fentmates.org\qquad\small\thepage}}
% la resta: el símbol a l'esquerra i el curs a la dreta, amb el títol al centre com abans
\fancyhead[L]{\raisebox{-0.3\height}{\includegraphics[height=0.62cm]{\fm@marca}}}
\fancyhead[R]{\small\bfseries\color{cursfosc}\fm@cursnom}
\renewcommand{\headrule}{\color{curs!45}\hrule height \headrulewidth}
\makeatother

\fmcurs{2batx-mii}

\begin{document}
\thispagestyle{primera}

\capcalera{Matemáticas II}{2\textsuperscript{o} Bach.}{Ficha de ejercicios · Unidad 8}{\raggedright Representación gráfica de funciones}
\vspace{10pt}
Nombre y apellidos: \rule{7.2cm}{0.4pt}\hfill Grupo: \rule{1.6cm}{0.4pt}\hfill Fecha: \rule{2.2cm}{0.4pt}
\vspace{8pt}

\begin{tcolorbox}[colback=soft,colframe=soft,arc=2mm,left=3mm,right=3mm,top=1.5mm,bottom=1.5mm]
\small Resuelve los ejercicios en el cuaderno justificando cada paso. Los ejercicios marcados con {\color{teal}\bfseries\faUniversity\ PAU} son de las pruebas de acceso a la universidad de la Comunidad Valenciana y de otras comunidades autónomas, traducidos al castellano. Puedes comprobar los resultados con las herramientas de los apuntes interactivos (dominios, asíntotas, estudio completo de una función paso a paso, transformaciones\dots).
\end{tcolorbox}

% =====================================================
\seccio{1}{Dominio de una función}

\ex{1} Calcula el dominio de las funciones siguientes:\begin{tasks}(2)
  \task $f(x)=\dfrac{x+3}{x^2-5x+6}$
  \task $f(x)=\sqrt{9-x^2}$
  \task $f(x)=\sqrt{\dfrac{x-1}{x+2}}$
  \task $f(x)=\ln\left(4-x^2\right)$
  \task $f(x)=\dfrac{1}{\sqrt{x^2-2x-3}}$
  \task $f(x)=\sqrt[3]{\dfrac{x}{x^2-1}}$
  \task $f(x)=\dfrac{e^x}{e^{2x}-1}$
  \task $f(x)=\ln(\ln x)$
\end{tasks}\recorda{Denominadores $\neq0$; radicandos de índice par $\geq0$ (y $>0$ si la raíz está en el denominador); argumentos de los logaritmos $>0$. El dominio es la intersección de todas las condiciones.}

\ex{2} Dominio de:\begin{tasks}(2)
  \task $f(x)=\tan3x$
  \task $f(x)=\dfrac{1}{1-\cos x}$
  \task $f(x)=\arcsin\dfrac x2$
  \task $f(x)=\sqrt{\sin x}$, en $[0,2\pi]$
  \task $f(x)=\ln(\cos x)$, en $[0,2\pi]$
\end{tasks}

\ex{2} Dominio de:\begin{tasks}(1)
  \task \pau{Aragón, ordinaria 2023} $f(x)=\dfrac{2x-1}{\sqrt{x^2-x-2}}$
  \task \pau{Canarias, ordinaria 2024} $f(x)=\dfrac{\ln(x+2)}{3x+4}$
\end{tasks}

\ex{3} Dominio de:\begin{tasks}(1)
  \task $f(x)=\sqrt{\dfrac{x^2-4}{x^2-9}}$
  \task $f(x)=\log_2\left(x^2-x\right)-\sqrt{3-x}$
\end{tasks}


% =====================================================
\seccio{2}{Cortes con los ejes y signo de la función}

\ex{1} Halla los puntos de corte con los ejes y estudia el signo de:\begin{tasks}(2)
  \task $f(x)=x^3-x^2-6x$
  \task $f(x)=\dfrac{x^2-4x+3}{x+1}$
  \task $f(x)=\left(x^2-1\right)e^x$
  \task $f(x)=\ln\left(x^2-8\right)$
\end{tasks}\recorda{Los ceros de $f$ y los puntos donde no es continua separan intervalos de signo constante: basta con un valor de prueba en cada uno.}

\ex{2} \pau{Extremadura, junio 2018} Demuestra que $f(x)=x^3+x-3$ corta el eje de abscisas en un único punto, que es positivo, y halla un intervalo de longitud $\frac12$ que lo contenga.

\ex{2} Estudia el signo de $f(x)=\dfrac{(x-2)^2(x+1)}{x^3}$. ¿Dónde corta o toca el eje OX?

\ex{3} ¿Cuántas soluciones tiene la ecuación $x=2\sin x$? Localiza la positiva entre dos décimas consecutivas.


% =====================================================
\seccio{3}{Simetrías y periodicidad}

\ex{1} Estudia la simetría de las funciones siguientes:\begin{tasks}(2)
  \task $f(x)=x^4-3x^2+1$
  \task $f(x)=\dfrac{x}{x^2-4}$
  \task $f(x)=e^{-x^2}\cos x$
  \task $f(x)=x^3+1$
  \task $f(x)=\ln\dfrac{1+x}{1-x}$
  \task $f(x)=\dfrac{\sin x\cos x}{x}$
\end{tasks}\recorda{Par: $f(-x)=f(x)$ (simetría respecto a OY). Impar: $f(-x)=-f(x)$ (respecto al origen). El dominio tiene que ser simétrico.}

\ex{1} Calcula el período de:\begin{tasks}(3)
  \task $\sin3x$
  \task $\cos\dfrac x2$
  \task $\sin x+\sin2x$
  \task $\sin^2(2x)$
  \task $\tan\dfrac{\pi x}{4}$
  \task $|\cos x|$
\end{tasks}

\ex{2} Comprueba que la gráfica de $f(x)=\dfrac{x^2+1}{x-1}$ es simétrica respecto al punto en el que se cortan sus asíntotas, y que la de $g(x)=\dfrac{1}{x^2-6x+10}$ es simétrica respecto a una recta vertical.

\ex{3} Demuestra que $f(x)=\ln\left(x+\sqrt{x^2+1}\right)$ es impar.


% =====================================================
\seccio{4}{Asíntotas verticales y horizontales}

\ex{1} Calcula las asíntotas verticales y horizontales (con los límites laterales):\begin{tasks}(2)
  \task $f(x)=\dfrac{2x^2-3}{x^2-1}$
  \task $f(x)=\dfrac{x+2}{x^2-4}$
  \task $f(x)=\dfrac{3e^x}{e^x+1}$
  \task $f(x)=\dfrac{\sqrt{x^2+1}}{x}$
  \task $f(x)=\dfrac{\ln x}{x-1}$
\end{tasks}\recorda{AV $x=a$ si algún límite lateral en $a$ es infinito. AH $y=k$ si $\displaystyle\lim_{x\to\pm\infty}f(x)=k$: estudia cada lado por separado. Cuidado: $\sqrt{x^2}=|x|$.}

\ex{2} \pau{Andalucía, ordinaria 2020} Estudia las asíntotas y la monotonía de $f(x)=\dfrac{x^2-2x-3}{x^2-1}$.

\ex{2} \pau{País Vasco, junio 2018} Asíntotas, monotonía y extremos de $f(x)=\dfrac{x^2}{x^2-4}$.

\ex{2} \pau{Cantabria, septiembre 2019} Dominio, asíntotas, monotonía y extremos de $f(x)=\dfrac{x+4}{x^2-7x-8}$.

\ex{2} \pau{País Vasco, extraordinaria 2025} Asíntotas y monotonía de $f(x)=\dfrac{x}{x^2-3x-4}$, y recta tangente en $x=0$.

\ex{3} \pau{La Rioja, junio 2018} Sea $f(x)=x\,e^{-ax}$. Estudia sus asíntotas según los valores de $a$ y halla $a$ para que tenga un extremo relativo en $x=1$. ¿Es máximo o mínimo?


% =====================================================
\seccio{5}{Asíntotas oblicuas y ramas parabólicas}

\ex{1} Calcula todas las asíntotas de:\begin{tasks}(1)
  \task \pau{Andalucía, junio 2017} $f(x)=\dfrac{x^2}{x-1}$ (halla también sus extremos).
  \task \pau{Asturias, extraordinaria 2020} $f(x)=\dfrac{2x^3+1}{x^2}$ (y los extremos).
\end{tasks}\recorda{$m=\displaystyle\lim_{x\to\pm\infty}\frac{f(x)}{x}$ (finito y $\neq0$), $n=\displaystyle\lim_{x\to\pm\infty}\big(f(x)-mx\big)$. En las racionales, la AO es el cociente de la división.}

\ex{2} \pau{Cantabria, ordinaria 2025} Sea $f(x)=\dfrac{x^3-8}{x^2-4}$ si $x\neq2$ y $f(2)=e^k$.\begin{tasks}(1)
  \task Halla $k$ para que sea continua en $x=2$.
  \task Calcula sus asíntotas.
  \task Escribe la recta tangente en $x=1$.
\end{tasks}

\ex{2} \pau{Canarias, ordinaria 2021} Sea $f(x)=\dfrac{ax^2-2}{b-x}$.\begin{tasks}(1)
  \task Halla $a$ y $b$ para que $f(-2)=2$ y $f$ sea continua en $\mathbb R-\{5\}$.
  \task Para $a=-1$, $b=-3$, calcula las asíntotas.
\end{tasks}

\ex{2} \pau{Castilla y León, julio 2018} Halla $a$ y $b$ para que $f(x)=\dfrac1x+ax+b$ tenga un mínimo relativo en $\left(\frac12,6\right)$ y, con estos valores, halla sus asíntotas.

\ex{3} \pau{La Rioja, julio 2018} Halla $a$ para que $\displaystyle\lim_{x\to+\infty}\left(\sqrt{9x^2+ax+1}-(3x-1)\right)=2$. ¿Cuál es, entonces, la asíntota oblicua de $y=\sqrt{9x^2+ax+1}$ en $+\infty$?

\ex{2} Estudia las ramas infinitas (asíntotas o ramas parabólicas) de:\begin{tasks}(2)
  \task $f(x)=\sqrt x$
  \task $f(x)=x+\sqrt x$
  \task $f(x)=x\ln x$
  \task $f(x)=\sqrt{x^2+x+1}-x$
\end{tasks}

\ex{2} Halla $b$ para que $f(x)=\dfrac{2x^2+bx}{x+1}$ tenga la asíntota oblicua $y=2x-1$.


% =====================================================
\seccio{6}{Monotonía, extremos y curvatura}

\ex{1} \pau{Cataluña, extraordinaria 2022} Sea $f(x)=\dfrac{9}{x^2+x-2}$. Determina su dominio, las asíntotas, los extremos relativos y los intervalos de crecimiento y decrecimiento.\recorda{Tabla de signos de $f'$ con los ceros de $f'$ y los puntos que no son del dominio; $f''$ para la curvatura.}

\ex{2} \pau{Extremadura, ordinaria 2021} Asíntotas, monotonía, extremos relativos y puntos de inflexión de $f(x)=e^{-x^2}$.

\ex{2} Estudia la monotonía, los extremos y la curvatura de $f(x)=x^4-4x^3$.

\ex{2} Estudia la monotonía y la curvatura de $f(x)=x^2\ln x$ y de $g(x)=\dfrac{x-1}{x^2}$.

\ex{3} Halla $a$, $b$ y $c$ para que $f(x)=x^3+ax^2+bx+c$ tenga un punto de inflexión en $(1,2)$ con tangente de pendiente $-3$.


% =====================================================
\seccio{7}{Representación de funciones polinómicas y racionales}

\ex{1} Haz el estudio completo y representa \pau{País Vasco, julio 2019} $f(x)=x^3+3x^2-2$.\recorda{Dominio, simetría, cortes, asíntotas (y posición), monotonía y extremos, curvatura e inflexión; después situamos toda la información y dibujamos.}

\ex{2} Haz el estudio completo y representa \pau{Asturias, extraordinaria 2024} $f(x)=\dfrac{x^2-4}{1-x}$.

\ex{2} Haz el estudio completo y representa \pau{La Rioja, extraordinaria 2020} $f(x)=\dfrac{x+3}{(x+2)^2}$ y calcula la recta tangente en el punto de inflexión.

\ex{2} Haz el estudio completo y representa \pau{Comunidad Valenciana, extraordinaria 2025} $f(x)=\dfrac{2x}{x^2+1}$ y calcula sus extremos absolutos.

\ex{3} Haz el estudio completo y representa $f(x)=\dfrac{x^3}{(x-1)^2}$.


% =====================================================
\seccio{8}{Representación de funciones radicales}

\ex{1} Dominio, cortes y asíntotas de:\begin{tasks}(2)
  \task $f(x)=\sqrt{x^2-9}$
  \task $f(x)=\dfrac{x+1}{\sqrt{x^2+3}}$
\end{tasks}\recorda{$\sqrt{x^2+bx+c}\approx|x|+\frac b2$ cuando $x\to\pm\infty$: las asíntotas pueden ser distintas en cada lado. Multiplica por el conjugado para calcular $n$.}

\ex{2} Haz el estudio completo y representa $f(x)=\dfrac{x}{\sqrt{x-1}}$.

\ex{2} Haz el estudio completo y representa $f(x)=x-3\sqrt[3]{x}$.

\ex{2} Haz el estudio completo y representa $f(x)=\dfrac{x+1}{\sqrt{x^2+3}}$.

\ex{3} Sea $f(x)=x\sqrt{4-x^2}$. Estudia su dominio, la simetría, los extremos absolutos y los puntos en los que la tangente es vertical, y represéntala.


% =====================================================
\seccio{9}{Representación de funciones exponenciales y logarítmicas}

\ex{1} Haz el estudio completo y representa \pau{Castilla y León, junio 2018} $f(x)=x\,e^{-x}$.\recorda{Dominio, simetría, cortes, asíntotas (y posición), monotonía y extremos, curvatura e inflexión; después situamos toda la información y dibujamos.}

\ex{2} Haz el estudio completo y representa \pau{País Vasco, extraordinaria 2024} $f(x)=2x\,e^{-2x^2}$.

\ex{2} Haz el estudio completo y representa \pau{Cantabria, extraordinaria 2022} $f(x)=\dfrac{e^x}{x}$.

\ex{2} Haz el estudio completo y representa \pau{Cataluña, ordinaria 2021} $f(x)=\dfrac{\ln x}{x}$.

\ex{2} Haz el estudio completo y representa $f(x)=\ln\left(x^2-4\right)$.

\ex{2} \pau{Castilla y León, ordinaria 2023} Dominio, monotonía y extremos de $f(x)=x(\ln x-1)$. ¿Cómo es la gráfica cerca de $x=0$?


% =====================================================
\seccio{10}{Representación de funciones trigonométricas}

\ex{2} Haz el estudio completo y representa $f(x)=\sin x\,(1+\cos x)$.\recorda{Si es periódica, estúdiala en un período. Factoriza $f'$ con las identidades ($\cos2x=2\cos^2x-1$…).}

\ex{2} Haz el estudio completo y representa $f(x)=\sin^2x+\cos x$ en $[0,2\pi]$.

\ex{2} Estudia $f(x)=\tan x-2x$ en $\left(-\frac\pi2,\frac\pi2\right)$: asíntotas, extremos y curvatura.

\ex{2} \pau{Cantabria, extraordinaria 2020} Sea $f(x)=\dfrac{1-\cos x}{x}$. Calcula $\displaystyle\lim_{x\to0}f(x)$ y las asíntotas de $f$.

\ex{2} Estudia la monotonía y la curvatura de $f(x)=x+\cos x$ en $[0,2\pi]$ y represéntala.


% =====================================================
\seccio{11}{Funciones a trozos, con valor absoluto y transformaciones}

\ex{1} Escribe a trozos y representa:\begin{tasks}(2)
  \task $f(x)=\left|x^2-2x-3\right|$
  \task $f(x)=x^2-4|x|+3$
  \task $f(x)=e^{-|x|}$
  \task $f(x)=|x|\,(x-2)$
\end{tasks}\recorda{$|g(x)|=g(x)$ si $g(x)\geq0$ y $-g(x)$ si $g(x)<0$. Los cambios de trozo suelen ser puntos angulosos.}

\ex{2} A partir de las gráficas de las funciones elementales, dibuja:\begin{tasks}(2)
  \task $y=\sqrt{x+4}-2$
  \task $y=-\sqrt{-x}$
  \task $y=2-e^{-x}$
  \task $y=\left|\ln(x+1)\right|$
\end{tasks}

\ex{2} \pau{Galicia, junio 2017} Calcula $a$ y $b$ para que sea derivable en $x=3$ la función $f(x)=\begin{cases}ax^2+b&\text{si }x<3\\ \ln(x-2)&\text{si }x\geq3\end{cases}$ y represéntala.

\ex{3} \pau{Madrid, extraordinaria 2022} Sea $f(x)=\begin{cases}x&\text{si }x\leq0\\ x\ln x&\text{si }x>0\end{cases}$.\begin{tasks}(1)
  \task Estudia la continuidad y la derivabilidad en $x=0$.
  \task Estudia su monotonía y sus extremos relativos, y represéntala.
\end{tasks}


% =====================================================
\seccio{12}{Información a partir de las gráficas de f, f' y f''}

\ex{2} Sabemos que $f'(x)=(x-1)^2(x+2)$. Sin calcular $f$, estudia su monotonía, los extremos y la curvatura, y dibuja una gráfica posible.\recorda{$f'>0\Rightarrow f$ crece; cambio de signo de $f'$ $\Rightarrow$ extremo; $f'$ crece ($f''>0$) $\Rightarrow f$ convexa; extremo de $f'$ $\Rightarrow$ inflexión de $f$.}

\ex{2} \pau{Cataluña, extraordinaria 2021} Calcula $a$ y $b$ para que $g(x)=ax^3+bx^2+1$ tenga un punto de inflexión en $x=\frac12$ y su derivada en ese punto sea $-\frac32$.

\ex{2} La gráfica de la derivada de un polinomio $p$ de grado 3 es una parábola que corta el eje OX en $x=-1$ y $x=2$. Si $p(0)=1$ y $p(2)=-19$, halla $p$ y di dónde tiene los extremos y el punto de inflexión.

\ex{2} \pau{Galicia, extraordinaria 2026} Dibuja la gráfica de una función $f:\mathbb R-\{2\}\to\mathbb R$ tal que $f$, $f'$ y $f''$ tengan el mismo signo en $(-1,2)$ y $\displaystyle\lim_{x\to2^+}f(x)=-\infty$. Explica el dibujo.

\ex{3} \pau{Galicia, ordinaria 2022} Dibuja la gráfica de una función continua y no negativa en $[0,3]$ tal que $f(0)=f(3)=0$, $f''>0$ en $(0,1)$, $f''<0$ en $(2,3)$ y constante en $(1,2)$. ¿Puede ser derivable en $x=1$?

\vfill
\begin{center}\color{muted}\small\faCheckCircle[regular]\ \ ¡Ánimo! Comprueba tus estudios con las herramientas de los apuntes interactivos.\end{center}

\end{document}
